\documentclass[11pt]{article}
\usepackage[margin=2.5cm]{geometry}
\usepackage[T1]{fontenc}
\usepackage[utf8]{inputenc}
\usepackage[slovene]{babel}
\usepackage{amsmath,amssymb,amsfonts}
\usepackage{array}
\usepackage{booktabs}
\usepackage{hyperref}
\usepackage{enumitem}

\begin{document}

\setlist[itemize]{leftmargin=1.5em, itemsep=0.25em, topsep=0.3em}
\renewcommand{\arraystretch}{1.2}

\section*{Rešitve in točkovnik: preizkus racionalnih funkcij}

\subsection*{1. naloga}

\[
 x^2-4=(x-2)(x+2),\qquad x^2-x-2=(x-2)(x+1).
\]
Imenovalec je ničeln pri $x=2$ in $x=-1$, zato
\[
 D_f=\mathbb{R}\setminus\{-1,2\}.
\]
Za dovoljene vrednosti lahko izraz krajšamo:
\[
 f(x)=\frac{x+2}{x+1},\qquad x\ne -1,2.
\]
Pri $x=2$ je bila prvotna funkcija izločena; poenostavljeni izraz bi tam imel vrednost $\frac{2+2}{2+1}=\frac{4}{3}$. Graf ima zato luknjo v točki $\left(2,\frac{4}{3}\right)$.

\vspace{0.5em}
\textbf{Točkovnik:}
\begin{itemize}[leftmargin=2.5em]
  \item razcep števca in imenovalca: 1 t,
  \item definicijsko območje: 2 t,
  \item poenostavitev: 1 t,
  \item luknja in koordinate: 1 t.
\end{itemize}

\subsection*{2. naloga}

Pogoj je $x\ne 1$. Za dovoljene vrednosti pomnožimo z $x-1$:
\[
 \frac{2x+1}{x-1}=3
 \iff 2x+1=3(x-1)
 \iff 2x+1=3x-3
 \iff x=4.
\]
Vrednost $4$ ni izločena, zato je rešitev $\boxed{x=4}$. Izločena vrednost je $x=1$.

\textbf{Točkovnik:}
\vspace{0.35em}
\begin{itemize}[leftmargin=2.5em]
  \item pogoj: 1 t,
  \item pravilno odstranjevanje imenovalca: 1 t,
  \item reševanje enačbe: 2 t,
  \item preverjanje in odgovor: 1 t.
\end{itemize}

\subsection*{3. naloga}

Kritični vrednosti sta $x=-3$ (imenovalec je nič, zato je vrednost izločena) in $x=2$ (števec je nič). Predznaki:

\begin{tabular}{|l|c|c|c|}
\hline
Interval & $(-\infty,-3)$ & $(-3,2)$ & $(2,\infty)$ \\
\hline
$\frac{x-2}{x+3}$ & $+$ & $-$ & $+$ \\
\hline
\end{tabular}

Iščemo negativne vrednosti, zato je rešitev
\[
 \boxed{(-3,2)}.
\]
Obe krajišči sta izključeni: pri $x=-3$ izraz ni definiran, pri $x=2$ pa je vrednost ulomka enaka 0, ne pa manjša od 0.

\textbf{Točkovnik:}
\vspace{0.35em}
\begin{itemize}[leftmargin=2.5em]
  \item kritične vrednosti: 2 t,
  \item pravilni predznaki: 2 t,
  \item izbira intervala: 1 t,
  \item pravilen zapis z izključenima krajiščema: 1 t.
\end{itemize}

\subsection*{4. naloga}

Najprej razcepimo števec:
\[
 3x^2+x-2=(3x-2)(x+1).
\]

\textbf{a)} Imenovalec ne sme biti nič, zato je $D_h=\mathbb{R}\setminus\{1\}$. (1 t)

\textbf{b)} Ničli sta $x=\frac{2}{3}$ in $x=-1$, torej točki $\left(\frac{2}{3},0\right)$ in $(-1,0)$. Za presečišče z osjo $y$ vstavimo $x=0$: $h(0)=2$, torej $(0,2)$. (3 t)

\textbf{c)} Ker se imenovalec pri $x=1$ izniči, števec pa tam ni ničeln, je navpična asimptota $x=1$. Deljenje polinomov da
\[
 h(x)=3x+4+\frac{2}{x-1},
\]
zato je poševna asimptota $y=3x+4$. (3 t)

\textbf{d)} Skica naj prikaže dve veji, navpično asimptoto $x=1$, poševno asimptoto $y=3x+4$ ter presečišči z osema. (1 t)

\subsection*{5. naloga}

Pogoja sta $x\ne 2$ in $x\ne -2$. Pomnožimo enačbo z $(x-2)(x+2)$:
\[
 \frac{1}{x-2}+\frac{1}{x+2}=1
 \iff (x+2)+(x-2)=x^2-4
\]
\[
 \iff 2x=x^2-4
 \iff x^2-2x-4=0.
\]
Po kvadratni formuli dobimo
\[
 x=\frac{2\pm\sqrt{4+16}}{2}=1\pm\sqrt{5}.
\]
Nobena vrednost ni $2$ ali $-2$, zato sta obe rešitvi dovoljeni:
\[
 \boxed{x=1-\sqrt{5}\quad\text{ali}\quad x=1+\sqrt{5}}.
\]

\textbf{Točkovnik:}
\vspace{0.35em}
\begin{itemize}[leftmargin=2.5em]
  \item pogoja: 1 t,
  \item pravilno odpravljanje ulomkov: 2 t,
  \item kvadratna enačba: 2 t,
  \item rešitvi: 1 t,
  \item preverjanje pogojev in končni odgovor: 1 t.
\end{itemize}

\subsection*{6. naloga}

Naj bo $x>0$ čas v urah, v katerem stroj B sam izdela naročilo. Stroj A v eni uri opravi $\frac{1}{6}$ naročila, stroja skupaj pa $\frac{1}{4}$ naročila. Zato velja
\[
 \frac{1}{6}+\frac{1}{x}=\frac{1}{4}.
\]
Odštejemo $\frac{1}{6}$:
\[
 \frac{1}{x}=\frac{1}{4}-\frac{1}{6}=\frac{3-2}{12}=\frac{1}{12},
\]
zato je $x=12$. Stroj B bi naročilo sam izdelal v \textbf{12 urah}.

\textbf{Točkovnik:}
\vspace{0.35em}
\begin{itemize}[leftmargin=2.5em]
  \item pomen neznanke in enota: 1 t,
  \item pravilna enačba: 3 t,
  \item pravilno odštevanje in računanje: 2 t,
  \item rešitev in preverjanje $x>0$: 1 t,
  \item odgovor z enoto: 1 t.
\end{itemize}


\end{document}
